\documentclass[11pt]{article}
\usepackage[margin=1in]{geometry}
\usepackage{amsmath,amssymb,amsthm}
\usepackage[T1]{fontenc}
\usepackage{lmodern}
\usepackage{microtype}
\usepackage[hidelinks]{hyperref}
\newtheorem{theorem}{Theorem}
\newtheorem{lemma}{Lemma}
\newcommand{\Pp}{\mathbb P}
\newcommand{\E}{\mathbb E}
\newcommand{\1}{\mathbf 1}
\newcommand{\Sc}{\mathcal S}
\newcommand{\Fc}{\mathcal F}
\newcommand{\Om}{\Omega}
\newcommand{\Rd}{\mathcal R_d}
\title{T21: Composition for the reported HPM\\\large C1 and discovery routes (a)--(f)}
\author{Astra lead; reviewed by Sol 6.1 x2}
\date{22 October 2026}
\begin{document}
\maketitle
\noindent\textbf{proved (Astra lead; reviewed by Sol 6.1 x2); not adopted.}
This proof establishes the conditional statement in \texttt{STATEMENT.md}.
The posterior concentration input H03, the addition input H10, and the terminal
repair input H24 are hypotheses, not previously proved theorems.
No numerical experiment is used.

\section{Setting and the composition claim}
All notation and hypotheses are as in \texttt{STATEMENT.md}. In particular,
\(\Om=\{S\subseteq[p]:|S|\le K\}\), \(K=\min(p,\lfloor n/2\rfloor)\),
\(S_0\in\Om\), and
\[
 F(S)=\log\{\varpi(S)m(S)\},\qquad
 Z=\sum_{S\in\Om}e^{F(S)},\qquad
 \pi(S)=e^{F(S)}/Z,\qquad \delta_1=1-\pi(S_0).
\]
The score formula and the capped prior are those written out in the statement.
The data law is \(P_0\); \(\Pp\) adds the independent algorithm randomness.
The total HPM order first compares scores, then (only for exact equality)
smaller size and lexicographic order. Its maximum on \(\Om\) is \(S^\dagger\).
The cumulative set of scored models is \(\Sc\), with pre-final subset
\(\Sc_{\rm cal}\). On \(U^c\), where a usable output exists, the reported
\(\widehat S\) is the maximum of \(\Sc\) in that same order.

The probabilities below are never conditioned on passing future diagnostics
or on future completion. In route (a), a launched final run executes all its
designated draws; in (d), the specified scan is completed; in (e), the specified
terminal rounds execute. These are the operational hypotheses in the statement.
If external interruption truncates a launched run, the full-count formula of
(a) is not asserted for that run.

\begin{theorem}[C1]
For every valid problem and every such joint law,
\begin{equation}\label{eq:C1}
 \Pp\{U\text{ or }\widehat S\ne S_0\}
 \le \Pp(U)+2E_0\delta_1+\Pp\{U^c,S_0\notin\Sc\}.
\end{equation}
Under H03, \(2E_0\delta_1\to0\). Consequently, if \(\Pp(U)\to0\)
and the discovery term tends to zero under one of the conditions established
below, \(\Pp\{U^c,\widehat S=S_0\}\to1\).
\end{theorem}

\begin{proof}
On \(\{\pi(S_0)>1/2\}\), every \(S\ne S_0\) satisfies
\[
 \pi(S)\le1-\pi(S_0)<\pi(S_0).
\]
Thus \(S_0\) is the strict score maximum, with no tie. If it has also been
scored and the output is usable, the method must return it. Hence
\[
 \{U^c,\widehat S\ne S_0\}
 \subseteq\{\delta_1\ge1/2\}\cup\{U^c,S_0\notin\Sc\}.
\]
The first event depends only on the data, and
\(\Pp(\delta_1\ge1/2)=P_0(\delta_1\ge1/2)\le2E_0\delta_1\).
Add \(\Pp(U)\) to get \eqref{eq:C1}.

H03 gives \(\delta_1\to0\) in \(P_0\)-probability, since
\(\sqrt{\log n/n}\to0\). For every \(\epsilon>0\), using
\(0\le\delta_1\le1\),
\[
 E_0\delta_1\le\epsilon+P_0(\delta_1>\epsilon).
\]
Taking a limsup, then \(\epsilon\downarrow0\), proves the asserted mean limit.
This does not assert an \(O(\sqrt{\log n/n})\) frequentist error rate.
\end{proof}

This argument also proves directly the ledger fact needed here: if the global
tie-broken maximum has been scored, the reported maximum on a usable run equals
it. No unproved clause of T11 or T03 is needed for this deterministic fact.

\section{Route (a): fresh final draws}
Let \(\mathcal H_{\rm cal}\) contain the data and the entire pre-final
history, including the tables and calibration streams. The event \(\Fc\)
that calibration permits the final SMC run, its tier, grid, \(T\), and
\(N_e\) are measurable there. On \(\Fc\), the number of initial draws
plus independence proposals is
\begin{equation}\label{eq:Q}
 Q_f=4N_e+4N_e(T+L+4)=4N_e(T+L+5).
\end{equation}
Off \(\Fc\) set \(Q_f=0\) in indicator-weighted formulas.

The wave-1 result T11(iii) says that, given data and frozen calibration, these
designated draws are iid \(\pi_0\), and their fixed-model miss probability is
\((1-\pi_0(S^*))^{Q_f}\). Its proof gives the same assertion conditional on
the entire pre-final history: see the collected proof's proposition on the
final run as a reading scheme, which fixes all calibration streams before
reading fresh final streams. For completeness, the relevant reasoning is as
follows. CAL and EX-S leave the final streams independent of
\(\mathcal H_{\rm cal}\). Each designated call uses unread uniforms and
has conditional law \(\pi_0\) given all preceding reads. Repeated conditioning
therefore factors the joint law of its \(Q_f\) outputs, even though other
random operations occur between calls. Their number is fixed by the pre-final
history, not selected by their outcomes. Thus, writing \(\mathcal D_f\) for
the set of these draws,
\begin{equation}\label{eq:miss}
 \Pp(S_0\notin\mathcal D_f\mid\mathcal H_{\rm cal})
 =(1-\pi_0(S_0))^{Q_f}\quad\hbox{on }\Fc.
\end{equation}
Here the true support is fixed by the data-generating problem and in particular
does not depend on final randomness. Initial draws and independence proposals
are all scored, so \(\Sc_{\rm cal}\cup\mathcal D_f\subseteq\Sc\) on
\(\Fc\). Splitting according to \(\Fc\), then applying
\eqref{eq:miss} and the tower property, gives
\begin{align}
 \Pp(S_0\notin\Sc)
 &\le\Pp(\Fc^c,S_0\notin\Sc_{\rm cal})\notag\\
 &\quad+\E\left[\1_{\Fc}\1_{\{S_0\notin\Sc_{\rm cal}\}}
       (1-\pi_0(S_0))^{Q_f}\right].\label{eq:a}
\end{align}
No independence between calibration success and its earlier scoring events
was used. The indicator of earlier absence is measurable before final
randomness. Since \(1-x\le e^{-x}\), the second line is at most the
same expectation with \(e^{-Q_f\pi_0(S_0)}\).

For any fixed \(M>0\),
\[
 \E[\1_{\Fc}e^{-Q_f\pi_0(S_0)}]
 \le\Pp(\Fc,Q_f\pi_0(S_0)\le M)+e^{-M}.
\]
The stated divergence on launched runs therefore makes this expectation
vanish, by first taking \(n\to\infty\) and then \(M\to\infty\).
Together with the nonlaunch-and-no-earlier-hit term, this proves (a).
This use of a fixed exponent does not apply to the adaptive calibration draw
count.

\subsection*{One-deletion start mass and actual budgets}
Assume \(s_0\ge1\), \(i\in S_0\), and
\(\widetilde S_0=S_0\setminus\{i\}\) has \(\pi_0\)-mass at least
\(c>0\). Its being the start mode is allowed but not needed for the following
inequality. The addition is feasible because \(s_0\le K\). Hypothesis H10
gives
\[
 \pi_0(S_0)\ge c r_{s_0-1}
 =\frac{c s_0}{(2p-s_0)\sqrt{n+1}}.
\]
Consequently, on \(\Fc\),
\begin{equation}\label{eq:intensity}
 Q_f\pi_0(S_0)\ge
 \frac{4c s_0N_e(T+L+5)}{(2p-s_0)\sqrt{n+1}}.
\end{equation}
Since \(p\ge s_0\ge1\), \(p\le2p-s_0\le2p\). For fixed
\(s_0,c\), the expression on the right has order
\(N_e(T+L+5)/(p\sqrt n)\), with constants independent of the tier and
grid. This is the order of a guaranteed lower bound, not a two-sided order
for the actual intensity. If \(p\le C\sqrt n\), then \(N_e\ge n\),
\(T\ge1\), and
\[
 \frac{4c s_0N_e(T+L+5)}{(2p-s_0)\sqrt{n+1}}
 \ge\frac{2c s_0}{C}\sqrt{\frac n{n+1}}(L+6)\longrightarrow\infty.
\]
This includes \(p\) proportional to \(\sqrt n\), since
\(L\ge\lceil\log_2(n+1)\rceil\). If the mass hypothesis holds only on
an event \(G_n\) with \(\Pp(G_n)\to1\), split the expectation in
\eqref{eq:a} over \(G_n\) and \(G_n^c\); the extra term is at most
\(\Pp(G_n^c)\). A lower bound which ceases to diverge gives no upper
bound on the true intensity and no proof of failure outside this sufficient
range. No conclusion of T27 is invoked.

\section{Route (b): a Gibbs execution}
At the specified execution put \(A=X_{{\rm pre}-j}\setminus\{j\}
=S_0\setminus\{j\}\). The two candidate configurations are \(A\) and
\(A\cup\{j\}\). If \(j\in S_0\), the latter is \(S_0\) and is
feasible because \(s_0\le K\). If \(j\notin S_0\), the former is
\(S_0\) and is feasible, even if the alternative exceeds the cap. In either
case SR scores \(S_0\). This is true whether or not the ensuing Gibbs draw
selects it. Thus a missed truth implies the absence of the specified
execution, and
\(\Pp(U^c,S_0\notin\Sc)\le\Pp(U^c,E_G^c)\).

The state immediately before coordinate \(j\) may have been changed by any
earlier labels in the random-order round. Therefore a condition solely on
round-start occupancy cannot be substituted for the displayed equality at
execution time. No round-end occupancy claim is being made.

\section{Route (c): stopped swap opportunities}
Suppose the current proxy has the form
\(S=(S_0\setminus\{l\})\cup\{j\}\), where
\(l\in S_0\), \(j\notin S_0\), and \(k=|S|=s_0\in\{1,\ldots,p-1\}\).
There is exactly one deletion/insertion pair producing \(S_0\), namely
\((j,l)\). Conditional on the full past just before a uniform swap,
that pair has probability
\begin{equation}\label{eq:swap}
 b=\frac1k\frac1{p-k}.
\end{equation}
Acceptance is irrelevant because the proposal is scored.

We prove the stopped-opportunity bound directly, rather than citing an
unproved ledger clause. Order all executed elementary operations
chronologically; the finite budgets permit such an ordering, and stopped
histories include the identity of the particle and the next operation.
Count only eligible swap executions before the first truth score. Let
\(\tau_i\) be the time of the \(i\)-th such execution, or \(\infty\)
if it does not occur. Eligibility and the absence of prior discovery are
known immediately before each draw. Let \(A_i\) be the event that
\(\tau_i<\infty\) and its proposal misses \(S_0\), and put
\(A_0\) equal to the whole probability space. The event
\(\{\tau_i<\infty\}\) implies \(A_{i-1}\) for \(i\ge2\): a hit at
an earlier counted execution would have stopped counting. Summing the
conditional miss probabilities over the possible pre-execution histories
gives
\[
 \Pp(A_i)=(1-b)\Pp(\tau_i<\infty)
 \le(1-b)\Pp(A_{i-1})\le(1-b)^i.
\]
This argument does not condition on a future opportunity count. If the
whole run misses truth and \(O\ge D\), then \(A_D\) occurs (for
\(D=0\), simply use the bound 1). Hence
\begin{align*}
 \Pp(S_0\notin\Sc)
 &\le\Pp(S_0\notin\Sc,O<D)
     +\Pp(S_0\notin\Sc,O\ge D)\\
 &\le\delta+(1-b)^D.
\end{align*}
The proof remains valid after conditioning on data or initial information
for which the stipulated proposal laws hold and \(D,b\) are fixed.
For PX choose \(D_n=\lceil\omega_n k(p-k)\rceil\). Then
\((1-b)^{D_n}\le e^{-bD_n}\le e^{-\omega_n}\to0\); the deficient-exposure
term vanishes by PX. A relation between two models, such as ONE, alone
does not supply this execution count.

For the neighbour-swap proposal, the unique useful pair is still
\((j,l)\). Its conditional probability is zero if
\(l\notin N(j)\setminus S\), and otherwise is
\[
 b_e=\frac1{k\,|N(j)\setminus S|}.
\]
The actual graph is essential. For a fixed deterministic \(b_*\in(0,1]\), if \(b_e\ge b_*\) at every counted
execution, the same conditioning argument yields
\(\Pp(A_i)\le(1-b_*)\Pp(A_{i-1})\), proving the asserted extension.
For a random lower bound measurable at a conditioning boundary, apply this argument conditionally and then average; the resulting unconditional term is the expectation of \((1-b_*)^D\). There is no division by zero: an execution with the required edge has a
nonempty insertion set. The uniform-swap formula was restricted to
\(1\le k<p\) precisely to exclude the method's self-loop boundaries.

\section{Route (d): tables and completed scans}
For \(n\ge4\), \(\lfloor\sqrt n\rfloor\le\lfloor n/2\rfloor\);
for \(n=3\), \(\lfloor\sqrt n\rfloor=1=\lfloor n/2\rfloor\).
Also \(B\le p\), so \(B\le K\). Thus every subset of a block is
feasible and every completed block table includes its true full-response
score. In particular, if \(S_0\subseteq C_b\) and that table is
constructed, \(S_0\) has been scored.

The feasible flip-or-swap neighbourhood is symmetric: the reverse of an
addition is a deletion, the reverse of a deletion is an addition to the
original feasible state, and a swap is reversed by the opposite swap.
If a completed scan has centre \(c\) in the neighbourhood of \(S_0\),
then \(S_0\) lies in the neighbourhood scanned at \(c\). The method
scores all of that neighbourhood, including at its last extra scan.
This is also the scoring portion of wave-1 T20(ii); it needs no strict-mode
assumption, although T20(ii)'s subsequent move/certificate conclusions do.
It follows that
\[
 \{U^c,S_0\notin\Sc\}\subseteq\{U^c,E_{TS}^c\}.
\]
Once a score has been obtained it remains in the HPM ledger even if the
schedule cannot be constructed later. Consequently schedule success is not
needed for this route, provided the output is usable. C1 already accounts
for output unavailability.

\section{Route (e): terminal repair under H24}
We prove the conditional repair inequality from the uniform one-round
bounds assumed in H24. Fix the data and the entire available past at the
selection boundary. On the availability event an eligible particle starts
at \(x\in\Rd\). Its future terminal rounds use fresh randomness and the
same kernel \(K_1\), without intervening resampling. The round count
\(H\), and the constants \(a=a_n\), \(\eta=\eta_n\), are now fixed.
H24 states
\begin{equation}\label{eq:H24}
 \inf_{x\in\Rd}K_1(x,\{S_0\})\ge a,
 \qquad
 \sup_{x\in\Rd}\Pp_x(\hbox{an exit during one round})\le\eta.
\end{equation}
These are assumptions on the actual hybrid round, including all its labels;
they are not derived here from leakage or gain conditions.

Let \(J\) be the event of any elementary-state exit from \(\Rd\) during
the \(H\) rounds. Before the first exit each round starts in \(\Rd\).
At each such boundary apply the second inequality in \eqref{eq:H24}
conditional on that boundary's past. Summing the probabilities of a first
exit in the different rounds gives
\begin{equation}\label{eq:exit}
 \Pp(J\mid\hbox{selection past})\le H\eta.
\end{equation}
For \(h=0,\ldots,H\), let \(B_h\) mean no exit through round \(h\)
and none of its first \(h\) round-end states equals \(S_0\).
With \(B_0\) the whole conditional space, on \(B_{h-1}\) the next start
is in \(\Rd\), so the first inequality in \eqref{eq:H24} gives
\[
 \Pp(B_h\mid\hbox{selection past})
 \le(1-a)\Pp(B_{h-1}\mid\hbox{selection past})\le(1-a)^h.
\]
Every terminal endpoint is a scored state. Indeed initial particle states
are scored; Gibbs updates score both feasible alternatives under SR, swap
and independence proposals are scored, and an acceptance moves to a scored
state while a rejection retains the current scored state. Resampling before
the terminal segment merely copies scored states. Thus a truth endpoint
implies a truth score. Let \(M_H\) be the event that \(S_0\notin\Sc\) when
those \(H\) rounds end. It follows that \(M_H\) is contained in
\(J\cup B_H\), and
\begin{equation}\label{eq:repair}
 \Pp(M_H\mid\hbox{selection past})
 \le(1-a_n)^H+H\eta_n.
\end{equation}
If the initial state is already truth, its score is already present and
\(M_H\) is empty. In particular this case cannot worsen any
whole-run bound. More generally, a truth visit before an endpoint only
improves discovery. At \(H=0\) use the convention \((1-a)^0=1\).

Let \(A\) be the past-measurable availability event. Split whole-run failure
over \(A^c\) and \(A\), use \eqref{eq:repair}, cap the conditional
bound by 1, and integrate. This gives
\[
 \Pp(S_0\notin\Sc)
 \le\Pp(A^c)+\E\left[\1_A
   \min\{1,(1-a_n)^H+H\eta_n\}\right].
\]
For a final-run segment use
\(A_F=\Fc\cap\{\hbox{eligible particle available}\}\).
Splitting first into nonlaunch, launch without an eligible particle, and
\(A_F\), improves this to
\begin{align}
 \Pp(S_0\notin\Sc)
 &\le\Pp(\Fc^c,S_0\notin\Sc_{\rm cal})
   +\Pp(\Fc,\hbox{no eligible particle})\notag\\
 &\quad+\E\left[\1_{A_F}
     \min\{1,(1-a_n)^H+H\eta_n\}\right].\label{eq:finalrepair}
\end{align}
Selection may depend on all states and randomness observed so far, since
\eqref{eq:H24} is uniform over eligible states. It may not inspect future
randomness. At a pilot terminal boundary the same argument uses the pilot's
forward filtration: unread uniforms are fresh there. Conditioning instead
on the grid eventually produced by that pilot would reveal future information
and is not part of this argument.

For arbitrary \(H\), \((1-a_n)^H\le e^{-Ha_n}\); thus the conditions
\(Ha_n\to\infty\) and \(H\eta_n\to0\) in probability on availability
make the capped expectation tend to zero, by the same bounded-variable
argument used for C1. For \(H=L+4\), H24's explicit asymptotic inputs give
this directly: when \(s_0,d\) are fixed and \(B\le16\),
\[
 c_n=\frac1{(s_0+d+B+2)!}\ge\frac1{(s_0+d+18)!}>0.
\]
On an event of probability tending to one,
\(a_n\ge c_n/2\), while \(L\to\infty\); hence \(Ha_n\to\infty\)
in probability, and \(H\eta_n\to0\) is the other assumed input.
The availability and nonlaunch terms must still vanish.
The proposed formulas for \(a_n,\eta_n\) in the statement only specify
which H24 input is intended; they do not prove H24 or its rates.

\section{Route (f): completed exact enumeration}
When
\begin{equation}\label{eq:exact}
 \sum_{k=0}^K\binom pk\le(5H_0+4N_0)p,
\end{equation}
the method chooses enumeration before constructing the SMC components.
On its completion, every feasible model has been scored, so \(\Sc=\Om\).
The running-best rule uses precisely the global tie order and therefore
returns \(\widehat S=S^\dagger\). Its accumulated evidence and inclusion
numerators are, under EX-A,
\[
 Z=\sum_{S\in\Om}\varpi(S)m(S),\qquad
 Z_j=\sum_{\substack{S\in\Om\\j\in S}}\varpi(S)m(S).
\]
Consequently the evidence is exact and the primary PIP is \(Z_j/Z\).
No Monte Carlo PIP or secondary-pooling assertion is used.
This is finite enumeration, not an application of an unproved asymptotic
Monte Carlo theorem. Wave-1 T09(i) establishes the exact start normaliser
and sampler when those components are used; it is compatible with the
method but unnecessary for these direct enumeration identities.

For fixed \(p\ge1\), the left side of \eqref{eq:exact} is at most
\(2^p\), whereas
\((5H_0+4N_0)p\ge4np\to\infty\). Thus the branch inequality holds
eventually (and eventually \(K=p\)). This proves eventual branch selection,
not completion under a fixed external allocation. Conditional on completion,
the discovery term in C1 is zero; H03 is still needed for truth consistency.

\section{Scope and source locators}
\small
The method used is \emph{METHODOLOGY\_v3.md}, version 3.0 of 19 October 2026:
\S3 for budgets; \S5 for the posterior; \S6.2 for exact enumeration;
\S6.3 for tables and the start; \S6.5 for move proposals;
\S\S6.8--6.9 for calibration and fresh final streams; and \S6.10 for the
all-scored HPM and local scans. The only cited proved wave-1 statements
are in \emph{COLLECTED\_PROOFS\_v2.tex}: the T11(iii) chapter, its
``final run is a reading scheme'' proposition and theorem; the T20 chapter,
clause (ii-A), first scoring implication; and the T09 chapter, part (i),
quoted only for the exact start component. All facts required beyond these
and the explicit H03/H10/H24 hypotheses were proved above from the method.

For statistical interpretation, G/R/SE/BM-C are the proposed primitive
conditions behind T03, not a proved route to H03 in this document. No
EXP or EXP' expansion is invoked and no no-switch conclusion is made;
therefore no cluster-separation implication is being presumed. The result
keeps operational failure separate from posterior error and discovery.
It proves neither evidence nor PIP accuracy on the SMC branch.
\end{document}




